\section{Future work}
\label{discussion:future-work}

The first direction addresses the single-node run-time bound directly. An indexed or co-located GeoParquet configuration was not tested, and a fairer comparison would add one. The benchmark wrote each tier in \texttt{partition\_key} order rather than along a space-filling curve, so within a file the per-row-group bounding boxes can overlap and the pruning the format allows is not fully realized (Section~\ref{system-architecture-and-implementation:data-layer:geoparquet-on-adls-gen2}). A spatial sort, such as a Hilbert or Z-order ordering, would tighten the row-group extents and raise the share of row groups and files a query can skip, without changing the format. A secondary spatial index over GeoParquet, or co-locating the data with the engine to remove the object-storage round-trip, is the second lever and would address the sub-second latency the deployed path paid (Section~\ref{discussion:reach-and-limits}). Any expected gain is confined to the containment and range patterns and to selective scans. A containment index does not order by distance without a bounding-box predicate, so it would not move the k-nearest-neighbor search.

A second direction locates the crossover the small-tier comparison only bounded. Within the 2--16 workers tested the broadcast join approached the PostGIS minimum without overtaking it, so the worker count at which distribution would overtake a tuned single node on a small join is an extrapolation rather than a measured point (Section~\ref{discussion:distribution-and-scaling}). Extending the worker sweep and varying the problem granularity would locate that crossover or bound it from above.

A third direction revisits the partitioned join. It completed only at the small tier and exhausted executor memory at the medium and large tiers, leaving open whether that failure reflects the strategy or its configuration (Section~\ref{discussion:reach-and-limits}). Tuning executor memory, the partition count, and the grid type would test whether the partitioned join completes where it failed here, separating a configuration limit from a strategy limit.